

Fertility is falling almost everywhere. In the United States, the period
fertility rate fell from about 2.1 children per woman in 2007 to 1.62 in
2023, and the world as a whole is likely below replacement for the first
time.\footnote{World Development Indicators \citep{worldBankWDI}, the
series in Figure~\ref{fig:quant-us-fertility}: 2.12 in 2007 and 1.62 in 2023.
On the world as a whole, see
\citet{fernandezVillaverdeNorrickTerraIncognita2026}.} The decline has become
a central concern for policymakers. Pinpointing its causes is
hard,\footnote{\citet{kearneyPuzzleFallingUS2022} find no economic or policy
factor that accounts for it. \citet{kearneyLevineWhyFertilityLow2025} argue
for ``shifting priorities'' across cohorts, a change in norms around work,
parenting, gender roles and leisure, shaped in part by new technologies such
as the smartphone and social media. \citet{doepkeEconomicsFertilityNew2023}
emphasize how easily parents can combine careers and family.} but housing and
its affordability have a role to play in what comes next, and housing is
often at the center of the conversation.\footnote{Singapore, for example, reserves up to 40\% of new
public flats for couples with children \citep{singaporeFPPS}. In the United
States, the 2025 budget law created a \$1,000 federal deposit for each child
born from 2025 to 2028 \citep{publicLaw119-21}. In a fourteen-country survey,
housing is the third most cited reason for having fewer children than
desired \citep{unfpaSWP2025}.} The logic is intuitive. Children need space,
and young households, who have little savings, find that space costly to pay
for.

Many of the housing policies proposed or enacted for families work through
credit.\footnote{Hungary writes off 10 million forints of a family's home loan
at each birth from the second child \citep{hungaryCSOKPlus}, and Poland
repays part of a family home loan at a second or third birth
\citep{polandFamilyRepayment}. South Korea lends at subsidized rates to
households with a newborn \citep{koreaNewbornLoan}. In 2024, the Harris
campaign proposed up to \$25,000 in down-payment assistance for first-time
buyers \citep{harrisCampaign2024}.} The reason is that family space usually has to be
bought. In the United States, 43\% of owned homes have seven or more rooms,
against 7.5\% of rented homes.\footnote{2023 American Housing Survey.} A
couple that wants a large home must usually buy one, and buying requires a
down payment. But a down payment is only part of what a child
costs in housing. A family also pays for its extra space every year its
children live at home, whether it rents or owns, and it keeps paying in years
when earnings fall. If the down payment is what holds young couples back,
easier credit could raise births. If the recurring cost of space is what
holds them back, credit may mainly change who owns a home and when. The
answer also changes over time, because births today become households a
generation later: as births decline, fewer families compete for homes, and a
policy that raises fertility raises the price of space for the next
generation of parents.

